\subsection{Design for Manufacture}
\subsubsection{3D-Printed Base Plates}

The top and bottom plates are 3D-printed in PLA at a 0.2mm layer height. Each plate is 110mm in diameter and 8mm thick with M4 screw holes and racetrack-shaped countersink pockets. The plates went through multiple print iterations to resolve a flatness issue where the bottom surface was slightly slanted, which affected the stability of the load cell. This was fixed by adjusting bed leveling and print orientation to ensure a flat bottom surface.

The second part of the plate design will feature an internal section to hide away components like the microcontroller PCB and battery. The goal is to get everything contained so that the load cell can bend without issue, while still hiding the rest of the internal components. 

\subsubsection{Load Cell Assembly}

The 75mm bar-type strain gauge load cell is sandwiched between the two plates using four M4 x 10mm bolts---two through the top plate into one end of the load cell, and two through the bottom plate into the other end. This creates a cantilever arrangement where weight applied to the top plate causes measurable strain on the load cell. The assembly requires only an M4 hex key and can be taken apart and reassembled in a few minutes.

\subsubsection{Electronics}

The electronics will be contained inside a separate section of the 3D Printed base. The top section contains the load cell. The bottom section will have all of the mounting for the PCB, battery, and amplifier board. The wires for the selected load cell will come down through the bottom mounted plate through a small hole, then be soldered onto the contacts for the amplifier board. Other wiring will be self-contained in the bottom section. 

